\documentclass{article}
\usepackage[T1]{fontenc}
\usepackage{lmodern}
\usepackage{graphicx}
\usepackage{amsmath,amssymb}
\usepackage[a4paper,margin=2cm]{geometry}
\usepackage[absolute,overlay]{textpos}
\setlength{\TPHorizModule}{1mm}
\setlength{\TPVertModule}{1mm}
\usepackage[indonesian]{babel}
\usepackage{hyperref}
\setlength{\emergencystretch}{2em}
% unit: Halaman 93

\begin{document}

% === Halaman 93 (python/native) ===

93

\#define ROTL(a,b) (((a) << (b)) | ((a) >> (32 - (b))))

\#define QR(a, b, c, d) ( 



\textbackslash{}


a += b,  d \textasciicircum{}= a,  d = ROTL(d,16), 

\textbackslash{}


c += d,  b \textasciicircum{}= c,  b = ROTL(b,12), 

\textbackslash{}


a += b,  d \textasciicircum{}= a,  d = ROTL(d, 8), 

\textbackslash{}


c += d,  b \textasciicircum{}= c,  b = ROTL(b, 7))

\#define ROUNDS 20

void chacha\_block(uint32\_t out[16], uint32\_t const in[16])

\{


int i;


uint32\_t x[16];


for (i = 0; i < 16; ++i) 



x[i] = in[i];


// 10 perulangan × 2 ronde/perulangan = 20 ronde


for (i = 0; i < ROUNDS; i += 2) \{



// Ronde ganjil



QR(x[0], x[4], x[ 8], x[12]); // kolom 0



QR(x[1], x[5], x[ 9], x[13]); // kolom 1



QR(x[2], x[6], x[10], x[14]); // kolom 2



QR(x[3], x[7], x[11], x[15]); // kolom 3



// Ronde genap



QR(x[0], x[5], x[10], x[15]); // diagonal 1 



QR(x[1], x[6], x[11], x[12]); // diagonal 2



QR(x[2], x[7], x[ 8], x[13]); // diagonal 3



QR(x[3], x[4], x[ 9], x[14]); // diagonal 4


\}


for (i = 0; i < 16; ++i)



out[i] = x[i] + in[i];

\}

Sumber: Wikipedia

\end{document}
